\Needspace{8\baselineskip}\originalchapter{电阻壁与环形压力模}\label{ch:wallballoon}
\originalsection{有限导体为何改变壁稳定}
理想近壁的感应电流永久屏蔽扰动; 有限电导壁的电流在扩散时间后衰减. 先算无壁势能 \ensuremath{\delta}W\ensuremath{\infty} 与理想壁位置 b 的 \ensuremath{\delta}Wb. 若两者均正则无需壁才能理想稳定, 两者均负则理想壁也救不了. 特殊情形是 \ensuremath{\delta}W\ensuremath{\infty}<0、\ensuremath{\delta}Wb>0: 快 Alfv\'en 时间稳定, 慢壁时间有 resistive wall mode. 简化薄壁模型增长率尺度 $\gamma\tau_w\sim-\delta W_\infty/\delta W_b$, 具体几何因子由真空匹配给出. 旋转、反馈和动理学阻尼可改变实验结果, 不能由静态理想墙结果单独推断.

\originalsection{轴对称位置模：磁通守恒给恢复力}
22.615 Lecture 21前半用一个理想导电的细电流环解释$n=0$位置模。它忽略压强和内部磁场的复杂形状，把环半径、竖直位置记为$R,Z$，小截面半径$a$固定。回路自感与链磁通为
\begin{equation}
L(R)=\mu_0R\left[\ln\frac{8R}{a}-2\right],\qquad
\mathcal F=LI-\Phi_{\rm ext}=\text{常数},\quad
\Phi_{\rm ext}=2\pi\int_0^R B_Z(R',Z)R'\dd R'.
\end{equation}
这里$\mathcal F$是整回路的磁通，后文$\psi$是每弧度极向磁通，二者不能混同。扰动时电流$I$随位置变化；将$U=LI^2/2$在固定$\mathcal F$下微分，利用真空外场$\nabla\cdot\mathbf B_{\rm ext}=0,\nabla\times\mathbf B_{\rm ext}=0$，给
\begin{equation}
F_Z=2\pi RI B_R,\qquad
F_R=\frac{I^2}{2}L'-2\pi RI B_Z .
\end{equation}
平衡要求$B_R=0$、$B_Z=IL'/(4\pi R)$。Lecture 21第3页两处把这里需要的$L'$印成$I'$；由磁通约束直接微分便能辨认。对中平面上下对称的外场，定义$n_B=-R(\partial_RB_Z)/B_Z$，再线性化，得
\begin{equation}
K_Z=\left.\frac{\partial F_Z}{\partial Z}\right|_0
=-\frac{I^2L'}{2R}n_B,\qquad
K_R=\left.\frac{\partial F_R}{\partial R}\right|_0
=\frac{I^2}{2}\left[L''-\frac{L'^2}{2L}
+(n_B-1)\frac{L'}R\right].
\end{equation}
负的$K$才是恢复力。通常$L'>0$，故竖直要求$n_B>0$；水平要求
$n_B<1-RL''/L'+RL'/(2L)$，在大对数$\ell=\ln(8R/a)$的领先近似趋于$n_B\lesssim3/2$。取$R=1$m、$a=.1$m、$I=10^6$A，保留自感全对数式得水平临界$n_B=1.41422$：$n_B=.5$两向恢复，$n_B=2$水平失稳，$n_B=-.5$竖直失稳。此为细线环与给定外场的局部结论，非任意形状tokamak的通用阈值。

\originalsection{环形曲率与 ballooning}
环形场的外侧为坏曲率, 局部压力驱动倾向把位移集中在外侧, 沿场磁线张力则惩罚集中. 高 n 采用 ballooning 表示, 把周期极向角展开成无限场线角 \ensuremath{\chi}, 将径向快速相位和局部剪切放入系数. 每个磁面的能量形如
\begin{equation}
 W(\Psi)=\int_{-\infty}^{\infty}\left[A(\chi)|X'|^2+C(\chi)|X|^2\right]\dd\chi,
 \qquad -(AX')'+CX=\lambda MX,
\end{equation}
A>0 为线弯曲, C含压力-曲率驱动和剪切, M>0 为惯性. 不能只保留外侧负C而舍掉整条线的张力. 典型 s--\ensuremath{\alpha} 模型中的 $s=rq'/q$, $\alpha\propto-q^2Rp'/B^2$ 给稳定区域, 常有第一与第二稳定区; 模型假设限于大长径比与特定剖面.

\originalsection{Mercier、Suydam 与边界渐近}
求 ballooning 方程前需检查 \ensuremath{\chi}\ensuremath{\to}\ensuremath{\pm}\ensuremath{\infty} 的允许渐近. 若局域 interchange/Mercier 条件违背, 试探函数已能给负能, 不应再把未衰减的无限角解作为有限能本征模. 若条件通过, 从一端衰减解积分, 用零点与 Sturm 比较判断压力驱动不稳定. Suydam 研究圆柱共振附近的径向局部结构, Mercier 在环形几何中推广; 后者不是把所有 R 简单换成 R0 的同一公式. 课程21--22的重点是这一表示与适用前提, 本册保留其逻辑, 不把一组特定装置的算图当普适证明.

\originalsection{Lecture 21的具体场线能量}
为让抽象$A,C,M$有可追查的内容，按原讲义定义轴对称坐标
$\mathbf B=(F/R)\mathbf e_\varphi+(\nabla\psi\times\mathbf e_\varphi)/R$，
$B_p=|\nabla\psi|/R$，$dV=2\pi J\,d\psi\,d\chi$，其中$\chi$为沿磁面的极向角、$J>0$为坐标Jacobian。令$\mathbf n=\nabla\psi/|\nabla\psi|$，$\mathbf b=\mathbf B/B$，
$\mathbf t=(B_\varphi/B)\mathbf b_p-(B_p/B)\mathbf e_\varphi$，曲率分解为$\boldsymbol\kappa=\kappa_n\mathbf n+\kappa_t\mathbf t$。磁场线局部绕行和安全因子为
\begin{equation}
\widehat q(\psi,\chi)=\frac{JB_\varphi}{R},\qquad
q(\psi)=\frac1{2\pi}\int_0^{2\pi}\widehat q\,d\chi .
\end{equation}
高环向数$n$的eikonal扰动沿场相位不变，先极小化平行压缩和最高阶磁压缩，令剩余幅度为$X(\chi)$，Lecture 21第10页得到（略去正的公共因子）
\begin{equation}
\begin{split}
W(\psi)=\int_{-\infty}^{\infty}J\dd\chi\Bigg[
&(k_n^2+k_t^2)\left|\frac{X'}{JB}\right|^2\\
&-\frac{2\mu_0RB_p}{B^2}\frac{dp}{d\psi}
(k_t^2\kappa_n-k_tk_n\kappa_t)|X|^2\Bigg],\\
k_t&=\frac{nB}{RB_p},\qquad
k_n=nRB_p\int_{\chi_0}^{\chi}
\frac{\partial\widehat q}{\partial\psi}(\psi,\chi')\dd\chi'.
\end{split}
\end{equation}
这里$X(\pm\infty)=0$且能量有限；$\chi_0$是定位相位，要在每个磁面上改变它检查。普通$2\pi$周期单函数不兼容剪切产生的径向相位，原讲义先在无限$\chi$上求衰减准模，再将其周期平移求和。这个公式来自高$n$领先阶，不能据此删去$n=1$外kink的平行电流项。

\originalsection{完整Mercier量与原件末因子的修正}
Lecture 22第3页写$D_M<1/4$为局域interchange稳定条件，但原印公式末尾是裸$R^2B_p^2/(JB^2)$，随$\chi$改变且带有量纲，不能作为磁面上的无量纲$D_M$。本卷按上一节实际$W$作独立推导，明确作为编者修正而非MIT发布的勘误。设
\begin{equation}
A_M=\frac{R^2B_p^2}{JB^2},\quad
\Lambda_M=F\left(\mu_0p'_\psi F
-\frac{R^2B_p^2}{J}\partial_\psi\widehat q\right),\quad
\langle Q\rangle_M=
\frac{\int_0^{2\pi}Q\,A_M^{-1}\dd\chi}
{\int_0^{2\pi}A_M^{-1}\dd\chi},
\quad p'_\psi=\frac{dp}{d\psi},\ q'_\psi=\frac{dq}{d\psi}.
\end{equation}
这是特定的加权磁面平均，不能换成体积平均。完成平方并在大$|\chi|$区间对周期快速变化的导数极小化，所得有效线弯曲系数为
$A_{\rm eff}=\langle A_M\rangle_M=2\pi/\int_0^{2\pi}A_M^{-1}\dd\chi$，而慢包络的能量为
$W_{\rm asym}\propto(q'_\psi)^2A_{\rm eff}
\int[\chi^2|Y'|^2-D_M|Y|^2]\dd\chi$。由此修正为
\begin{equation}
\boxed{
D_M=\frac{\mu_0p'_\psi}
{(q'_\psi)^2\langle A_M\rangle_M^2}
\left[
2\left\langle\frac{RB_p\kappa_n}{B^2}\right\rangle_M
+\left\langle\frac{\Lambda_M}{B^4}\right\rangle_M
-\left\langle\frac1{B^2}\right\rangle_M
 \left\langle\frac{\Lambda_M}{B^2}\right\rangle_M
\right].}
\end{equation}
量纲可核：$[A_M]=\mathrm{T\,m}$，修正末因子是$(\mathrm{T\,m})^2$，使$D_M$无量纲且只依赖$\psi$。具体推导在本卷QA的独立核查记录中保留了完成平方、周期均匀化每一步。慢包络满足$(\chi^2Y')'+D_MY=0$，所以$Y\sim|\chi|^{-1/2\pm\sqrt{1/4-D_M}}$。$D_M>1/4$出现可给负能的局域振荡尾；$D_M<1/4$仅排除此局域机制，后续还须检查整条场线的ballooning方程。等号是临界。

圆柱极限直接检验修正：取$R=R_0$、$B_p=B_\theta$、$J=r/B_\theta$、$q=rB_z/(R_0B_\theta)$，磁面系数不再随$\chi$变，$\Lambda_M$的协方差消去，得
\begin{equation}
D_M=-\frac{2\mu_0p_r'q^2}{rB_z^2q_r'^2}=D_s,\qquad
D_M<\frac14\ \Longleftrightarrow\
rB_z^2(q_r'/q)^2+8\mu_0p_r'>0 .
\end{equation}
这与本卷上一章从Lecture 18第8、11页推得的Suydam式一致。一个局部、同时满足径向力平衡的例子：$R_0=3$m、$r=.5$m、$B_z=2$T、$B_\theta=.3$T、$q_r'=.8\,\mathrm{m^{-1}}$；$p_r'=-10^5$Pa/m时$D_M=.24240684<1/4$，改为$-1.2\times10^5$Pa/m得$.29088821>1/4$。这是单个磁面局部判别，不推断全局稳定。原讲义近轴非圆式还用未在这两讲定义中心归一化的$\beta_{p0}$；本稿不据该符号编造数值阈值。

\originalsection{可计算的圆截面 s-alpha 模型}
Lecture 22第5--6页另取大长径比圆截面、窄压强梯度层，$\chi\to\theta$、$J\to r/B_\theta$、$R\to R_0$、$B_p\to B_\theta$、$\kappa_n\to-\cos\theta/R_0$、$\kappa_t\to\sin\theta/R_0$。从$d\psi/dr=R_0B_\theta$可知$d p/d\psi=(R_0B_\theta)^{-1}dp/dr$；原件第5页把$R_0$误印为$r$，且第5--6页局部绕行分子将$JB_\varphi$误印为$JB_\theta$，均以Lecture 21定义和圆柱极限核查后在此修正。令
\begin{equation}
s=\frac r q\frac{dq}{dr},\qquad
\alpha=-\frac{2\mu_0r^2p_r'}{R_0B_\theta^2}
=-q^2R_0\frac{d\beta}{dr},\quad
\beta=\frac{2\mu_0p}{B_0^2},\qquad
\Lambda_B=s(\theta-\theta_0)-\alpha(\sin\theta-\sin\theta_0).
\end{equation}
外侧$\theta=0$坏曲率、向外压强减小$p_r'<0$给$\alpha>0$；$\theta_0$是需改变的相位。将上节具体$W$在这些近似下除以正系数，得到
\begin{equation}
W_{s\alpha}[X]=\int_{-\infty}^{\infty}
\left[(1+\Lambda_B^2)|X'|^2
-\alpha(\Lambda_B\sin\theta+\cos\theta)|X|^2\right]\dd\theta ,
\end{equation}
\begin{equation}
\boxed{\frac{d}{d\theta}\left[(1+\Lambda_B^2)X'\right]
+\alpha(\Lambda_B\sin\theta+\cos\theta)X=0},
\qquad X(-\infty)=X(+\infty)=0 .
\end{equation}
这是零频边界方程；若求实际增长率还须补惯性权重。$\Lambda_B'=s-\alpha\cos\theta$说明压强还改变局部剪切，故增加$\alpha$不必单调恶化稳定性。一个明确的允许试探位移取$\theta_0=0$、$X=\exp(-\theta^2/2)$，Gaussian积分给
\begin{equation}
\frac{W_{s\alpha}}{\sqrt\pi}
=\frac12+\frac34s^2-\alpha\e^{-1/4}\left(1+\frac74s\right)
+\frac{3-\e^{-1}}4\alpha^2.
\end{equation}
$s=1$、$\alpha=1.4$时$W_{s\alpha}\simeq-0.812925<0$，足以证明这一模型该点有负能方向；$\alpha=4$时同一试函数为正，不能据一个试函数宣布第二稳定区。真正判全域须遍历允许$X$、$\theta_0$与磁面，并检查前述Mercier边界条件。
